% Derived from upstream styles/warm-editorial/template-169.tex at e63b61b140a00596e5260ec7611e42bfde7efb5f.
\documentclass[aspectratio=169,10pt]{beamer}
\usepackage{fontspec}
\defaultfontfeatures[TeX Gyre Pagella]{Path=fonts/tex-gyre/,UprightFont=texgyrepagella-regular.otf,BoldFont=texgyrepagella-bold.otf,ItalicFont=texgyrepagella-italic.otf,BoldItalicFont=texgyrepagella-bolditalic.otf}
\defaultfontfeatures[TeX Gyre Heros]{Path=fonts/tex-gyre/,UprightFont=texgyreheros-regular.otf,BoldFont=texgyreheros-bold.otf,ItalicFont=texgyreheros-italic.otf,BoldItalicFont=texgyreheros-bolditalic.otf}

\usepackage{fontspec}
\usepackage{graphicx}
\usepackage{tikz}
\usepackage{amsmath}
\usepackage{unicode-math}
\usetikzlibrary{arrows.meta,backgrounds,calc,fit,positioning,shapes.geometric}

\setmainfont{TeX Gyre Pagella}
\setsansfont{TeX Gyre Heros}
\setmathfont[Path=fonts/tex-gyre/]{texgyrepagella-math.otf}

\definecolor{BMEIvory}{HTML}{F7F2E8}
\definecolor{BMEInk}{HTML}{2F3437}
\definecolor{BMEBrown}{HTML}{5C433B}
\definecolor{BMETerracotta}{HTML}{B66A50}
\definecolor{BMESage}{HTML}{7D9276}
\definecolor{BMESand}{HTML}{E7DAC7}
\definecolor{BMEStone}{HTML}{C9BBA8}
\definecolor{BMEMuted}{HTML}{77716B}
\definecolor{BMEWhite}{HTML}{FFFDFC}
\definecolor{BMEFlow}{HTML}{655A54}
\definecolor{BMEReturn}{HTML}{718C80}
\definecolor{BMEOptional}{HTML}{938B83}

\setbeamercolor{normal text}{fg=BMEInk,bg=BMEIvory}
\setbeamercolor{structure}{fg=BMEBrown}
\setbeamertemplate{navigation symbols}{}
\setbeamertemplate{footline}{}
\setbeamertemplate{headline}{}
\AtBeginEnvironment{frame}{\raggedright\hyphenpenalty=10000\exhyphenpenalty=10000\relax}

\newcommand{\WarmFrameHeader}[3]{%
  \node[anchor=north west,inner sep=0pt,text=BMETerracotta,
        font=\sffamily\fontsize{6.2}{7.5}\selectfont\bfseries]
    at ([xshift=0.86cm,yshift=-0.30cm]current page.north west) {\MakeUppercase{#2}};
  \node[anchor=north west,inner sep=0pt,text=BMEBrown,
        font=\rmfamily\fontsize{16.2}{19.2}\selectfont\bfseries]
    at ([xshift=0.86cm,yshift=-0.58cm]current page.north west)
    {\makebox[11.9cm][l]{#1}};
  \draw[BMETerracotta,line width=1.25pt]
    ([xshift=0.86cm,yshift=-1.29cm]current page.north west) -- ++(0.74cm,0);
  \draw[BMEStone,line width=0.45pt]
    ([xshift=1.72cm,yshift=-1.29cm]current page.north west) --
    ([xshift=-0.86cm,yshift=-1.29cm]current page.north east);
  \node[anchor=north east,inner sep=0pt,text=BMEBrown,
        font=\sffamily\fontsize{8.5}{10}\selectfont\bfseries]
    at ([xshift=-0.86cm,yshift=-0.58cm]current page.north east) {#3};
}

\newcommand{\WarmFrameFooter}[1]{%
  \node[anchor=south west,inner sep=0pt,text=BMEMuted,
        font=\sffamily\fontsize{5.8}{7}\selectfont]
    at ([xshift=0.86cm,yshift=0.20cm]current page.south west) {#1};
  \fill[BMETerracotta]
    ([xshift=-1.02cm,yshift=0.24cm]current page.south east) circle[radius=1.2pt];
}

\newcommand{\WarmTag}[2][BMESand]{%
  \tikz[baseline=(warmtag.base)]{%
    \node[rounded corners=3pt,fill=#1,draw=none,minimum width=0pt,minimum height=0pt,
          outer sep=0pt,inner xsep=3.6pt,inner ysep=1.2pt,text height=1.15ex,text depth=0.20ex,
          text=BMEBrown,font=\sffamily\fontsize{4.9}{5.8}\selectfont\bfseries] (warmtag) {#2};}%
}

\newcommand{\WarmBodyText}{%
  \color{BMEInk}\sffamily\mdseries\fontsize{6.8}{8.8}\selectfont
  \raggedright\rightskip=0pt plus 1fil\spaceskip=0.30em plus 0.05em minus 0.03em\xspaceskip=0.45em\relax}
\newcommand{\WarmBodySmall}{%
  \color{BMEInk}\sffamily\mdseries\fontsize{6.5}{8.3}\selectfont
  \raggedright\rightskip=0pt plus 1fil\spaceskip=0.30em plus 0.05em minus 0.03em\xspaceskip=0.45em\relax}
\newcommand{\WarmDisplayMath}{%
  \color{BMEBrown}\rmfamily\fontsize{13.5}{16.2}\selectfont}

\tikzset{
  slide/prose/.style={align=left,execute at begin node={\raggedright\hyphenpenalty=10000\exhyphenpenalty=10000\relax}},
  warm/panel/.style={rounded corners=3.5pt,draw=BMEStone,line width=0.45pt,fill=BMEWhite},
  warm/node/.style={rounded corners=2.8pt,draw=BMEStone,line width=0.55pt,fill=BMEWhite,
    minimum height=0.72cm,inner xsep=0.18cm,inner ysep=0.12cm,align=center,
    text=BMEInk,font=\sffamily\fontsize{6.5}{7.8}\selectfont},
  warm/action/.style={warm/node,align=left,inner xsep=0.22cm,inner ysep=0.16cm},
  warm/decision/.style={diamond,aspect=1.55,draw=BMEBrown,line width=0.62pt,fill=BMEWhite,
    minimum width=1.18cm,minimum height=0.76cm,inner sep=0.05cm,align=center,
    text=BMEBrown,font=\sffamily\fontsize{5.8}{6.9}\selectfont\bfseries},
  warm/terminal/.style={rounded corners=8pt,draw=BMEStone,line width=0.55pt,fill=BMEWhite,
    minimum height=0.58cm,inner xsep=0.22cm,inner ysep=0.10cm,align=center,
    text=BMEBrown,font=\sffamily\fontsize{6.0}{7.2}\selectfont\bfseries},
  warm/graph edge/.style={draw=BMEStone,line width=1.08pt,line cap=round,line join=round},
  warm/graph node/.style={circle,draw=BMEBrown,line width=0.75pt,fill=BMEWhite,
    minimum size=0.68cm,inner sep=0pt,text=BMEBrown,font=\sffamily\fontsize{9}{10}\selectfont\bfseries},
  warm/graph current/.style={warm/graph node,draw=BMETerracotta,line width=0.95pt,fill=BMESand},
  warm/graph route/.style={draw=BMETerracotta,line width=1.25pt,line cap=round,line join=round,
    shorten <=3.4mm,shorten >=3.4mm},
  warm/graph route arrow/.style={warm/graph route,-{Stealth[length=1.35mm,width=1.0mm]}},
  warm/flow/.style={-{Stealth[length=1.35mm,width=1.0mm]},draw=BMEFlow,line width=0.66pt,
    line cap=round,line join=round},
  warm/return/.style={-{Stealth[length=1.30mm,width=0.95mm]},draw=BMEReturn,line width=0.61pt,
    line cap=round,line join=round},
  warm/optional/.style={-{Stealth[length=1.25mm,width=0.92mm]},draw=BMEOptional,line width=0.56pt,
    dashed,dash pattern=on 3.1pt off 2.5pt,line cap=round,line join=round},
  warm/highlight/.style={-{Stealth[length=1.35mm,width=1.0mm]},draw=BMETerracotta,line width=0.78pt,
    line cap=round,line join=round},
  warm/arrow/.style={warm/flow},
  warm/edge label/.style={text=BMEMuted,font=\sffamily\fontsize{5.4}{6.6}\selectfont,fill=BMEIvory,
    rounded corners=1.5pt,inner xsep=2.4pt,inner ysep=1.1pt},
  warm/adjacent label/.style={text=BMEMuted,font=\sffamily\fontsize{5.4}{6.6}\selectfont,inner sep=0pt},
  warm/label/.style={warm/edge label}
}

\usepackage{xeCJK}
\setCJKmainfont[Path=fonts/shipporimincho/,BoldFont=ShipporiMincho-Bold.ttf]{ShipporiMincho-Regular.ttf}
\setCJKsansfont[Path=fonts/zenkakugothicnew/,BoldFont=ZenKakuGothicNew-Bold.ttf]{ZenKakuGothicNew-Regular.ttf}
% Japanese type adapter and user-requested ink navy; original macro geometry remains.
\definecolor{BMEInk}{HTML}{243746}
\definecolor{BMEBrown}{HTML}{243746}
\newcommand{\Base}[3]{\fill[BMEIvory] (current page.south west) rectangle (current page.north east);\WarmFrameHeader{#1}{#2}{#3}}
\newcommand{\Foot}{\WarmFrameFooter{学習用の仮定・試算 / 一つの窓口・受付順・一定処理時間 / 実測値ではありません}}
\tikzset{jptext/.style={anchor=north west,inner sep=0pt,align=left,font=\sffamily\fontsize{9}{13}\selectfont,text=BMEInk},jpsmall/.style={font=\sffamily\fontsize{7}{10}\selectfont,text=BMEMuted}}
\begin{document}

% Cover geometry follows the original Warm Editorial starter.
\begin{frame}[plain]
\begin{tikzpicture}[remember picture,overlay]
\fill[BMEIvory] (current page.south west) rectangle (current page.north east);
\fill[BMETerracotta] ([xshift=.88cm,yshift=-1.22cm]current page.north west) rectangle ++(.14cm,-.86cm);
\node[anchor=north west,inner sep=0pt,text=BMETerracotta,font=\sffamily\fontsize{7}{9}\selectfont\bfseries] at ([xshift=1.24cm,yshift=-1.25cm]current page.north west) {SYSTEMS / 学習ノート};
\node[anchor=north west,inner sep=0pt,text=BMEBrown,font=\rmfamily\fontsize{29}{37}\selectfont\bfseries,text width=8cm] at ([xshift=1.22cm,yshift=-2.60cm]current page.north west) {待ち時間の\\しくみ};
\node[jptext,text width=8cm,text=BMEMuted] at ([xshift=1.25cm,yshift=-5.15cm]current page.north west) {仕事の速さと、届く間隔。\\二つの時間から、列の長さを考える。};
\draw[BMEStone,line width=.55pt] ([xshift=1.25cm,yshift=-6.45cm]current page.north west) -- ++(7.3cm,0);
\node[anchor=south west,inner sep=0pt,jpsmall] at ([xshift=1.25cm,yshift=.72cm]current page.south west) {仮定に基づく5枚の教材 / 2026.10};
\begin{scope}[shift={([xshift=11.2cm,yshift=-3.1cm]current page.north west)}]
\foreach \y/\qname in {0/qone,-.9/qtwo,-1.8/qthree} {\node[warm/node,minimum width=2.5cm,minimum height=.65cm] (\qname) at (0,\y) {待っている仕事};}
\node[warm/node,draw=BMETerracotta,fill=BMESand,minimum width=2.5cm] (service) at (0,-3.25) {窓口で処理};
\draw[warm/highlight] (qthree.south) -- (service.north);
\node[jpsmall] at (0,.75) {順番が来るまで待つ};
\end{scope}
\end{tikzpicture}
\end{frame}

\begin{frame}[plain]\begin{tikzpicture}[remember picture,overlay]
\Base{一つの窓口に、仕事が届く}{MODEL / 前提}{02}
\node[jptext,text width=13cm] at ([xshift=.98cm,yshift=-1.72cm]current page.north west) {窓口は一つ。先に届いた仕事から、1件ずつ処理する。};
\begin{scope}[shift={([xshift=.98cm,yshift=-3.25cm]current page.north west)}]
\node[warm/terminal,minimum width=2.3cm,font=\sffamily\fontsize{9}{12}\selectfont] (a) at (1.2,0) {到着};
\node[warm/node,minimum width=2.3cm,font=\sffamily\fontsize{9}{12}\selectfont] (w) at (5,0) {順番待ち};
\node[warm/node,fill=BMESand,draw=BMETerracotta,minimum width=2.3cm,font=\sffamily\fontsize{9}{12}\selectfont] (p) at (8.8,0) {処理};
\node[warm/terminal,minimum width=2.3cm,font=\sffamily\fontsize{9}{12}\selectfont] (f) at (12.6,0) {完了};
\draw[warm/flow] (a.east) -- (w.west);\draw[warm/flow] (w.east) -- (p.west);\draw[warm/flow] (p.east) -- (f.west);
\end{scope}
\draw[BMEStone,line width=.45pt] ([xshift=.98cm,yshift=-4.65cm]current page.north west) -- ++(13.8cm,0);
\node[jptext,text width=5.8cm] at ([xshift=1.02cm,yshift=-5.05cm]current page.north west) {\color{BMETerracotta}\bfseries 処理時間 $p=5$ 秒\par\vspace{.25cm}\WarmBodyText 1件の仕事が、窓口を使う時間。\\途中で割り込まず、毎回同じ長さ。};
\node[jptext,text width=6cm] at ([xshift=8cm,yshift=-5.05cm]current page.north west) {\color{BMESage}\bfseries 到着間隔 $d$ 秒\par\vspace{.25cm}\WarmBodyText 次の仕事が届くまでの時間。\\この教材では一定の間隔にする。};
\Foot\end{tikzpicture}\end{frame}

\begin{frame}[plain]\begin{tikzpicture}[remember picture,overlay]
\Base{待ち時間を式にする}{RECURRENCE / 計算}{03}
% Heading and each formula are separate nodes; \strut gives every formula the same WarmDisplayMath line box.
\node[jptext,text width=6cm] at ([xshift=.98cm,yshift=-1.8cm]current page.north west) {\bfseries 前の仕事が終わったら、始める};
\node[jptext] at ([xshift=.98cm,yshift=-2.45cm]current page.north west) {\WarmDisplayMath\strut$s_i=\max(a_i,f_{i-1})$};
\node[jptext] at ([xshift=.98cm,yshift=-3.20cm]current page.north west) {\WarmDisplayMath\strut$f_i=s_i+p$};
\node[jptext] at ([xshift=.98cm,yshift=-3.95cm]current page.north west) {\WarmDisplayMath\strut$w_i=s_i-a_i$};
\node[jptext,text width=6.2cm] at ([xshift=8cm,yshift=-1.85cm]current page.north west) {\WarmBodyText $a_i$：到着時刻　$s_i$：開始時刻\\$f_i$：完了時刻　$w_i$：待ち時間\par\vspace{.45cm}最初の仕事は時刻0に到着。\\初期値は $f_0=0$ とする。};
\draw[BMEStone,line width=.45pt] ([xshift=.98cm,yshift=-4.95cm]current page.north west) -- ++(13.8cm,0);
\node[jptext,text width=13cm] at ([xshift=1cm,yshift=-5.25cm]current page.north west) {例：3件目は8秒に到着。前の仕事は10秒に完了。};
\begin{scope}[shift={([xshift=2cm,yshift=-6.85cm]current page.north west)}]
\draw[warm/flow] (0,0) -- (10.8,0);
\foreach \x/\t in {0/8,3/10,10.5/15} {\draw[BMEInk] (\x,-.1) -- (\x,.1);\node[jpsmall,below] at (\x,-.2) {\t 秒};}
\draw[BMESage,line width=3pt] (0,.3) -- (3,.3);\draw[BMETerracotta,line width=3pt] (3,.3) -- (10.5,.3);
\node[jpsmall,above] at (1.5,.5) {待ち $2$ 秒};\node[jpsmall,above] at (6.75,.5) {処理 $5$ 秒};
\end{scope}
\Foot\end{tikzpicture}\end{frame}

\begin{frame}[plain]\begin{tikzpicture}[remember picture,overlay]
\Base{間隔が1秒違うと、待ちは積み上がる}{COMPARISON / 仮の試算}{04}
\begin{scope}[shift={([xshift=1.45cm,yshift=-6.4cm]current page.north west)},x=1cm,y=.24cm]
\foreach \y in {0,5,10,15} {\draw[BMEStone,line width=.35pt] (0,\y) -- (7.4,\y);\node[jpsmall,left] at (-.1,\y) {\y};}
\draw[warm/flow] (0,0) -- (7.9,0);\draw[warm/flow] (0,0) -- (0,17.3);
\foreach \x in {0,...,7} {\pgfmathtruncatemacro{\n}{\x+1}\node[jpsmall,below] at (\x,-.3) {\n};}
\draw[BMETerracotta,line width=1.3pt] (0,0) -- (7,14);
\draw[BMEInk,line width=1.1pt] (0,0) -- (7,7);
\draw[BMESage,line width=1.2pt] (0,0) -- (7,0);
\foreach \x in {0,...,7} {\fill[BMETerracotta] (\x,2*\x) circle[radius=1.3pt];\fill[BMEInk] (\x,\x) circle[radius=1.3pt];}
\node[jpsmall,anchor=west,text=BMETerracotta] at (7.2,14) {$d=3$ 秒};\node[jpsmall,anchor=west,text=BMEInk] at (7.2,7) {$d=4$ 秒};\node[jpsmall,anchor=west,text=BMESage] at (7.2,1) {$d=6$ 秒};
\node[jpsmall,anchor=west] at (0,18) {待ち時間（秒）};\node[jpsmall,anchor=west] at (6,-2.5) {仕事の番号};
\end{scope}
\node[jptext,text width=4cm] at ([xshift=11cm,yshift=-2.15cm]current page.north west) {\bfseries 処理は毎回5秒\par\vspace{.45cm}\WarmBodyText 3秒間隔なら、\\1件進むごとに2秒ずつ\\待ちが増える。\par\vspace{.35cm}6秒間隔なら、この仮定では\\順番待ちは生じない。\par\vspace{.5cm}\WarmDisplayMath\raggedright $w_i=(i-1)$\par\vspace{.08cm}$\phantom{w_i=}\cdot\max(0,p-d)$};
\Foot\end{tikzpicture}\end{frame}

\begin{frame}[plain]\begin{tikzpicture}[remember picture,overlay]
\Base{待ち時間を減らす、三つの方法}{DESIGN / 考え方}{05}
\node[anchor=north west,inner sep=0pt,slide/prose,text width=3.8cm] at ([xshift=.98cm,yshift=-1.85cm]current page.north west) {\color{BMETerracotta}\sffamily\fontsize{7}{9}\selectfont\bfseries 01 / 処理を短く\par\vspace{.3cm}\color{BMEBrown}\rmfamily\fontsize{15}{19}\selectfont\bfseries 5秒を4秒に\par\vspace{.4cm}\WarmBodyText 4秒間隔の到着に間に合う。\\同じ仮定なら待ちは0秒。};
\node[anchor=north west,inner sep=0pt,slide/prose,text width=3.8cm] at ([xshift=5.55cm,yshift=-1.85cm]current page.north west) {\color{BMESage}\sffamily\fontsize{7}{9}\selectfont\bfseries 02 / 到着を分散\par\vspace{.3cm}\color{BMEBrown}\rmfamily\fontsize{15}{19}\selectfont\bfseries 4秒を5秒に\par\vspace{.4cm}\WarmBodyText 5秒の処理と到着間隔が揃う。\\変動があれば余裕も必要。};
\node[anchor=north west,inner sep=0pt,slide/prose,text width=3.8cm] at ([xshift=10.12cm,yshift=-1.85cm]current page.north west) {\color{BMETerracotta}\sffamily\fontsize{7}{9}\selectfont\bfseries 03 / 窓口を増やす\par\vspace{.3cm}\color{BMEBrown}\rmfamily\fontsize{15}{19}\selectfont\bfseries 同時に処理\par\vspace{.4cm}\WarmBodyText 仕事を複数の窓口に配分する。\\窓口ごとの空き時刻を計算。};
\draw[BMEStone,line width=.45pt] ([xshift=.98cm,yshift=-4.65cm]current page.north west) -- ++(13.8cm,0);
\node[jptext,text width=13cm] at ([xshift=1.02cm,yshift=-5.15cm]current page.north west) {\color{BMEBrown}\rmfamily\fontsize{16}{22}\selectfont\bfseries 平均だけでなく、ばらつきも見る\par\vspace{.4cm}\WarmBodyText 実際の仕事は、到着も処理も一定ではない。\\この式で関係を理解したら、次は到着時刻と処理時間の記録で確かめる。};
\Foot\end{tikzpicture}\end{frame}
\end{document}
